\scanpage{171}
\mspage{120}
% archive sheet number: 171
\uline{Corollaire 2} : Sous les conditions du cor.~1, supposons que $I$ ait un objet final $i_0$. Alors

\[
(\unsure{p_{i_0}|p\vert_{i_0}})_!(e_{X_{i_0}})\longrightarrow p_!(e_X)
\]
est un isomorphisme dans $Y^\wedge$.

D'où aussitôt la proposition 11 : si $X_{i_0}$ est \deleted{$Y_{/p_{i_0\,!}(e_{X_{i_0}})}$} tot. $W$-asph., alors \inserted{$X$ l'est, car} pour tout $p:X\to Y$, $p_!(e_X)\simeq p_{i_0\,!}(X_{i_0})$, et ce dernier élément de $Y^\wedge$ est $W$-asphérique par hyp. sur $X_{i_0}$.

Notons le cas particulier

\uline{Corollaire 3} Soit $Z$ \unsure{un|une} \deleted{catégorie}\inserted{objet} totalement $W$-asphérique dans $\mathrm{Cat}$. Alors son produit $Z\times I$ par \unsure{tte|tt} $I$ dans $\mathrm{Cat}$ ayant un objet final est aussi \inserted{tot.} $W$-asphérique.

\deleted{\uline{Corollaire \unsure{4}} Mais dans le cas d'une catégorie produit $X=Z\times I$, \unsure{considérée} comme cofibrée donc propre sur $I$, on peut faire mieux. Dans ce cas, le foncteur}

\[
\deleted{\unsure{i|j}\longmapsto p_{i\,!}(X_i)}
\]
